% =====================================================================
%  FICHE 4 — THEME 1 — Corrige niveau MOYEN
% =====================================================================
\documentclass[11pt,a4paper]{article}
\usepackage[utf8]{inputenc}
\usepackage[T1]{fontenc}
\usepackage{lmodern}
\usepackage[french]{babel}
\usepackage{amsmath,amssymb,mathtools}
\usepackage{icomma}
\usepackage[version=4]{mhchem}
\usepackage{siunitx}
\sisetup{output-decimal-marker={,},per-mode=symbol}
\DeclareSIUnit\Molar{M}
\usepackage{xcolor}
\usepackage{tikz}
\usepackage[most]{tcolorbox}
\usepackage{enumitem}
\usepackage{booktabs}
\usepackage{array}
\usepackage{fancyhdr}
\usepackage[a4paper,margin=2.2cm,headheight=26pt,headsep=10pt]{geometry}

\definecolor{primary}{RGB}{146,95,160}
\definecolor{primarydark}{RGB}{92,52,104}
\definecolor{excol}{RGB}{0,135,90}

\tcbset{boxbase/.style={enhanced,breakable,boxrule=0.9pt,arc=2.5pt,
   left=3mm,right=3mm,top=2mm,bottom=2mm,fonttitle=\bfseries,coltitle=white,
   toptitle=0.5mm,bottomtitle=0.5mm}}

\newcounter{exo}
\newenvironment{corr}[1][]{%
  \par\medskip\refstepcounter{exo}%
  \noindent\begin{tcolorbox}[boxbase,colback=excol!5,colframe=excol,
     title={\textsc{Correction \theexo}\ifx\relax#1\relax\relax\else\textmd{ --- \itshape #1}\fi}]}%
  {\end{tcolorbox}}

\pagestyle{fancy}\fancyhf{}
\fancyhead[L]{\small\textcolor{primarydark}{\textbf{Fiche 4} --- Th.\,1 : Mole, masse molaire, entit\'es}}
\fancyhead[R]{\small\textcolor{primarydark}{\textsc{Chimie} --- \textbf{Corrigé Moyen}}}
\fancyfoot[C]{\small\thepage}
\renewcommand{\headrulewidth}{0.8pt}
\renewcommand{\headrule}{\hbox to\headwidth{\color{primary}\leaders\hrule height \headrulewidth\hfill}}
\newcommand{\NA}{N_{\!\text{A}}}
\setlength{\parindent}{0pt}

\begin{document}

\begin{center}
{\LARGE\bfseries\color{primarydark} Thème 1 --- Corrigé Moyen}
\end{center}

\begin{corr}[chaîne masse $\to$ moles $\to$ molécules]
\textbf{a)} $M(\ce{C9H8O4})=9(12)+8(1)+4(16)=108+8+64=\SI{180}{\gram\per\mole}.$\\[2pt]
\textbf{b)} On convertit d'abord $\SI{360}{\milli\gram}=\SI{0.360}{\gram}$, puis
$m\xrightarrow{\div M}n\xrightarrow{\times\NA}N$ :
\[ n=\frac{m}{M}=\frac{0{,}360}{180}=\SI{2.00e-3}{\mole},\qquad
   N=n\times\NA=\num{2.00e-3}\times\num{6.02e23}=\num{1.204e21}. \]
\[ \boxed{N\approx\num{1.20e21}\ \text{molécules d'aspirine}} \]
\end{corr}

\begin{corr}[compter des ions à partir d'une masse]
$M(\ce{FeSO4})=56+32+4(16)=\SI{152}{\gram\per\mole}$. Avec $m=\SI{304}{\milli\gram}=\SI{0.304}{\gram}$ :
\[ n(\ce{FeSO4})=\frac{0{,}304}{152}=\SI{2.00e-3}{\mole}. \]
Chaque unité \ce{FeSO4} donne \textbf{un} ion \ce{Fe^2+}, donc $n(\ce{Fe^2+})=\SI{2.00e-3}{\mole}$ :
\[ N(\ce{Fe^2+})=n\times\NA=\num{2.00e-3}\times\num{6.02e23}=\num{1.204e21}. \]
\[ \boxed{N(\ce{Fe^2+})\approx\num{1.20e21}\ \text{ions}} \]
\end{corr}

\begin{corr}[pourcentage massique d'un élément]
\textbf{a)} $M(\ce{CH4N2O})=12+4+28+16=\SI{60}{\gram\per\mole}$. Le carbone représente $\SI{12}{\gram}$
sur $\SI{60}{\gram}$ par mole :
\[ \%\,\ce{C}=\frac{12}{60}\times100=\SI{20.0}{\percent}. \]
\textbf{b)} $M(\ce{NaHSO3})=23+1+32+3(16)=\SI{104}{\gram\per\mole}$. L'oxygène représente
$3\times16=\SI{48}{\gram}$ par mole :
\[ \%\,\ce{O}=\frac{48}{104}\times100\approx\SI{46.2}{\percent}. \]
\textit{Astuce de contrôle :} la somme des pourcentages de tous les éléments doit faire $100\,\%$.
\end{corr}

\begin{corr}[doses de micronutriments et conversions d'unités]
\textbf{a)} $\SI{14}{\milli\gram}=\SI{1.4e-2}{\gram}$, et $M(\ce{Fe})=\SI{56}{\gram\per\mole}$ :
\[ n(\ce{Fe})=\frac{\num{1.4e-2}}{56}=\SI{2.5e-4}{\mole}. \]
\textbf{b)} $M(\ce{C10H16N2O3S})=10(12)+16(1)+2(14)+3(16)+32=120+16+28+48+32=\SI{244}{\gram\per\mole}$.\\
Avec $\SI{48.8}{\micro\gram}=\SI{48.8e-6}{\gram}$ :
\[ n=\frac{\num{48.8e-6}}{244}=\SI{2.0e-7}{\mole}. \]
\[ \boxed{n(\ce{Fe})=\SI{2.5e-4}{\mole};\quad n(\text{B8})=\SI{2.0e-7}{\mole}} \]
\textit{Point de vigilance :} bien convertir mg et \si{\micro\gram} en grammes avant de diviser par $M$.
\end{corr}

\begin{corr}[comparer des nombres d'atomes]
On calcule $n$ pour chaque échantillon, puis le nombre d'\emph{atomes} (attention : \ce{O2} est
\textbf{diatomique}).
\[ \ce{He}:\ n=\frac{2{,}0}{4}=\SI{0.50}{\mole}\ \text{d'atomes}\ \Rightarrow\
   N=0{,}50\,\NA=\num{3.01e23}\ \text{atomes}. \]
\[ \ce{O2}:\ n=\frac{8{,}0}{32}=\SI{0.25}{\mole}\ \text{de molécules}\ \Rightarrow\
   N_{\text{atomes}}=2\times0{,}25\,\NA=0{,}50\,\NA=\num{3.01e23}\ \text{atomes}. \]
\[ \ce{Al}:\ n=\frac{9{,}0}{27}=\SI{0.333}{\mole}\ \text{d'atomes}\ \Rightarrow\
   N=0{,}333\,\NA=\num{2.0e23}\ \text{atomes}. \]
\[ \boxed{\ce{He}\ \text{et}\ \ce{O2}\ \text{en contiennent le plus}\ (\num{3.01e23}\ \text{atomes chacun})} \]
\textit{Leçon :} pour \ce{O2}, il faut \textbf{doubler} le nombre de molécules pour obtenir le nombre
d'atomes --- l'erreur serait de s'arrêter à $0{,}25$ mole.
\end{corr}

\end{document}
